% Section 5. Experiments.
% Every number carries a "% source:" comment naming a file under
% /home/miria/OptEvolve/v2/results/. Nothing here comes from atlas/.
% Style contract: no em-dash or en-dash characters, no triple-hyphen dash
% macro, no italic markup of any kind, pure ASCII, plain declarative prose.
% All figures below were recomputed from the raw JSON rows on 2026-09-10 by
% opening the cited file and agree with it to the digits printed. Where a
% recomputation disagreed with an earlier draft the earlier value was deleted,
% not softened.

\section{Experiments}
\label{sec:results}
\label{sec:exp}

All measurements run on an RTX 4090 (sm\_89) and an RTX PRO 6000 Blackwell
Max-Q (sm\_120), over four families: total variation denoising of six retinal
OCT scans, min-cost flow on grid and regular topologies, capacitated
multi-commodity flow, and optimal transport. Termination is decided by a
float64 relative duality gap in original coordinates at $10^{-4}$ unless stated
otherwise. Every time is wall clock to a certified gap and every iteration
count is the count at which that certificate passed. Iteration count is written
$N$ and per-iteration cost $t$, and we state which of the two any ratio is
measured on, because the thesis of the paper is $T = N \times t$.

\subsection{Problem families}
\label{sec:families}

% source: row and cell counts enumerated from every JSON under results/ carrying a
% per-instance 'cell' or 'family' key, plus the tolerance-keyed factorials for TV.
Table~\ref{tab:families} lists the five families on which every number in this
paper was measured, with the volume of measurement behind each. The three
linear-programming families share the same standard-form composite and are
solved by the same code path, which is what makes the transfer experiment of
Section~\ref{sec:exp:transfer} meaningful: the solver is held fixed and only the
operator changes. Total variation is the one non-LP family and the only one to
which the fused kernel of Section~\ref{sec:exp:fusion} applies, since that kernel is
a stencil.

\begin{table}[t]
\caption{Problem families with measurements. Rows count individual solves
recorded under \texttt{results/}. Cells count distinct problem configurations.}
\label{tab:families}
\begin{center}
\begin{tabular}{llrr}
\toprule
Family & Configurations & Rows & Cells \\
\midrule
Min-cost flow, grid & $k \times k$ grid, bidirectional 4-neighbour arcs & 5{,}645 & 4 \\
Min-cost flow, random regular & degrees 2, 3, 4; $n = 256, 512, 1024$ & 3{,}929 & 8 \\
Optimal transport & bipartite, $m = n \in \{32, 64, 96\}$ & 1{,}431 & 3 \\
Total-variation denoising & OCT scans, 128 to 1984 pixels & 1{,}003 & 7 sizes \\
Multi-commodity flow & capacity-coupled, $K \in \{2, 4\}$ & 648 & 2 \\
\bottomrule
\end{tabular}
\end{center}
\end{table}

The families are not variations on a single operator. Table~\ref{tab:opnorm}
gives the behaviour of $\|L\|$, which is the quantity the construction gate of
Section~\ref{sec:exp:method} tests and therefore the quantity that decides which
recipes are admissible at all. Two families have an operator norm bounded
independently of problem size, one grows without bound in the size, and one
contracts in its coupling parameter. This spread is the reason the admissible
region is family dependent rather than a property of the solver.

\begin{table}[t]
\caption{Operator-norm behaviour across the five families, measured. The
reported $\|L\|$ is the gate's bound, which carries a $1.02$ safety factor over
the spectral norm.}
\label{tab:opnorm}
\begin{center}
\begin{tabular}{lll}
\toprule
Family & $\|L\|$ behaviour & Measured \\
\midrule
% source: build_instances('grid', k) norm_bound; 528 to 9024 variables
Min-cost flow, grid & saturates at $4$, bounded in size & $4.045$ to $4.076$ over a $17\times$ size increase \\
% source: build_instances('regular', n, degree=d) norm_bound
Min-cost flow, regular & bounded by degree, $\leq \sqrt{4d}$ & $2.884$ at $d=2$, $3.935$ at $d=4$ \\
% source: atlas/bench/ot.py transport_composite norm_bound vs 1.02*sqrt(m+n)
Optimal transport & $\sqrt{m+n}$, unbounded in size & exact to $3 \times 10^{-14}$ relative \\
% source: atlas/bench/mcf.py mcf_composite norm_bound, grid size 8
Multi-commodity flow & grows with $K$, admissible $\tau\sigma$ contracts & $\|L\|^2$ from $17.1$ at $K{=}1$ to $150.1$ at $K{=}128$ \\
Total-variation denoising & gradient stencil, bounded & fixed by the discretisation \\
\bottomrule
\end{tabular}
\end{center}
\end{table}

The consequence is direct. On grid instances the admissible region is
effectively fixed: $\|L\|$ moves $0.76$ percent while the variable count grows
$17\times$, so the same step sizes remain legal at every size. On optimal
transport the admissible product $\tau\sigma$ must shrink as $1/(m+n)$, and both
catalog step-size settings become inadmissible at sizes far below those we
solve. On multi-commodity flow the admissible product falls from $0.0584$ at
$K=1$ to $0.00666$ at $K=128$. A step-size pair is therefore not portable across
families, and the gate detects this symbolically, before any iteration runs.

\subsection{The factorisation}
\label{sec:exp:factorisation}

% source: results/tv_ffi_20260906/cadence_4090_v1/factorial.json
% source: results/bw_20260909/sz256_bw/factorial.json
Table~\ref{tab:Nvectors} gives the per-instance iteration counts on TV at
$n=256$. The six-element vectors, not only their means, are identical element by
element on the two architectures at every cadence and under both schemes.

\begin{table}[t]
\centering
\caption{Per-instance iteration counts over the six OCT scans at $n=256$,
tolerance $10^{-4}$. Every vector is identical on the RTX 4090 and on the
Blackwell. The vectors are also identical across the float64, float32-iterate
and two-phase precision policies at $K=64$ and $K=256$. At $K=16$ they are not:
the float32-iterate policy moves PDHG from the float64 vector
$[1008, 496, 608, 400, 320, 640]$, mean $578.7$, to
$[1008, 512, 608, 416, 320, 640]$, mean $584.0$, while Condat-Vu is unchanged
at $584.0$.}
\label{tab:Nvectors}
\begin{tabular}{@{}llrr@{}}
\toprule
Cadence & Scheme & Per-instance $N$ & Mean \\
\midrule
$K=16$  & Condat-Vu & 1008, 512, 608, 416, 320, 640 & 584.0 \\
$K=16$  & PDHG      & 1008, 496, 608, 400, 320, 640 & 578.7 \\
$K=64$  & both      & 1024, 512, 640, 448, 320, 640 & 597.3 \\
$K=256$ & both      & 1024, 512, 768, 512, 512, 768 & 682.7 \\
\bottomrule
\end{tabular}
\end{table}

% source: results/tv_ffi_20260906/{cadence_4090_v1,f32row_4090,f64host_4090,kernelhost_4090}/factorial.json
% source: results/bw_20260909/{sz128_bw,sz496_bw}/factorial.json
% The width counterexamples are the reason this paragraph is scoped. They were
% recomputed from the two bw_20260909 files and from f32row_4090 named above.
The same vectors appear in four separate factorial artifacts covering the
float64 device loop, the float32-iterate float64-certificate device loop, the
host-phase policy and the hand written kernel in both execution structures. At
$n=256$ and $K \geq 64$ iteration count is invariant to arithmetic width,
execution structure, kernel and device, and is a function
of cadence and tolerance alone. Width is not inert in general. On the Blackwell
at $K=64$ float64 gives $576.0$ against float32's $586.7$ at $n=128$ and
$1066.7$ against $1077.3$ at $n=496$, and on the 4090 at $K=16$ float32 moves
PDHG from $578.7$ to $584.0$. Cadence quantises the width difference away at
$K \geq 64$ and $n=256$ and not everywhere.

% source: results/kernel_fusion_20260909/e2e/full_4090_sz{128,256,384,496}.json
% source: results/kernel_fusion_20260909/e2e_bw/full_bw_sz{128,256,384,496}.json
% Within-cell iteration identity and the 6-of-6 certification were rechecked
% arm by arm; the largest within-cell time spread over the eight cells is 2.6294.
The invariance survives into the deployed solver. Across eight arms, the traced
baseline, the static loop and six fusion tiles, at four sizes on both devices,
the per-instance vectors are identical within each cell and every arm certifies
6 of 6, while solve time within a cell varies by up to $2.629\times$. The sweep
requested seven tiles, $8\times4$, $16\times2$, $16\times4$, $24\times4$,
$32\times2$, $32\times3$ and $32\times4$. The $32\times3$ tile never runs,
because the harness requires the fused depth $T$ to divide the cadence $K$ and
no recorded cadence, $64$, $128$, $256$ or $512$, is divisible by 3. Six tiles
are measured everywhere in this paper.

% source: results/netflow_20260909/grid_4090.json
% 48 rows: 24 banked (all executed) and 24 balanced (all gate-rejected).
% The two rows at grid size 12, instance 1, hit the 20000 cap at gap 2.858e-4.
The second family agrees. All 24 banked min-cost flow rows return exactly equal
iteration counts under the static and the traced loop, and 22 of the 24
certify; the two rows at grid size 12, instance 1, agree only because both hit
the 20000-iteration cap at a gap of $2.858 \times 10^{-4}$. Over the 22
certified rows the traced-over-static wall-time ratio has geometric mean
$1.395$ and spans $1.293$ to $1.492$; over all 24, including the two capped
rows, it is $1.385$ and spans $1.277$ to $1.492$. All 12 matched instance pairs
return exactly equal counts under PDHG and Condat-Vu, of which 11 are certified
pairs and one is the pair that both schemes ran to the cap.

% source: results/method_axis_20260909/factorial.json (XLA half)
% source: results/tv_ffi_20260906/kernelhost_4090/factorial.json (kernel half)
% Both halves recomputed as solver_s / iters, averaged over the six images.
Per-iteration cost is invariant to scheme identity under XLA and is not under
the hand written kernel. At \texttt{f64\_K256} under XLA the two schemes return
the identical vector $[1024, 512, 768, 512, 512, 768]$, penalised SGM-1 times in
ratio $0.99494$ and per-iteration costs of $19.76$ against $19.76$
microseconds. In the custom-kernel artifact, with the same
\texttt{build\_digest} in both rows, Condat-Vu costs $19.21$ microseconds per
iteration and PDHG $25.02$, a factor of $1.303$, and the gap appears in all
four custom-kernel cells at $1.303$, $1.264$, $1.179$ and $1.197$. Scheme
identity is inert to $t$ on the compiler path and not on the kernel path.

% source: results/tv_ffi_20260906/cadence_4090_v1/factorial.json and
% results/bw_20260909/sz256_bw/factorial.json and
% results/method_axis_20260909/factorial.json
% CORRECTION relative to results/tv_ffi_20260906/DECOMPOSITION_REPORT.md, which
% states 584/597/683 for both schemes. That is true at K=64 and K=256 only.
Scheme identity is exactly inert to $N$ at coarse cadence and not at fine
cadence. At $K=16$ the schemes differ on two of the six instances, 512 against
496 and 416 against 400, and the split reproduces on the Blackwell. At
$10^{-6}$ and $K=64$ Condat-Vu averages $2860.0$ iterations and PDHG $2870.7$;
both certify 5 of 6, both means include the same capped instance, and the two
schemes differ on exactly one of the five certified instances, 2432 against
2496. A coarse cadence quantises a real scheme difference away, which is the
mechanism this paper argues for rather than an exception to it.

\subsection{Execution structure}
\label{sec:exp:execution}

% source: results/artifacts_20260910/fanova.json, block
% exec_x_prec_x_cadence__logT__tol0.0001, which carries all twelve cells to
% full precision and names its four source factorials:
%   device_loop|f64  -> tv_ffi_20260906/cadence_4090_v1/factorial.json
%   device_loop|f32  -> tv_ffi_20260906/f32row_4090/factorial.json
%   host_phases|f64  -> tv_ffi_20260906/f64host_4090/factorial.json
%   host_phases|f32  -> tv_ffi_20260906/px999999_4090/factorial.json
% The twelve values are 22.536573, 20.389612, 19.777049, 22.986976, 20.405764,
% 19.775866, 27.184947, 16.185603, 13.290695, 27.998956, 12.228629, 8.431614.
Table~\ref{tab:periter} decomposes per-iteration cost over execution structure,
arithmetic width and cadence.

\begin{table}[t]
\centering
\caption{Microseconds per iteration on the RTX 4090, TV at $n=256$, Condat-Vu,
averaged over the six scans. PDHG in the same artifacts gives $27.40$, $17.20$
and $14.15$ at host-phase float64 and $26.59$, $12.21$ and $8.36$ at host-phase
float32. The host-phase float32 row is the patched policy registry and not the
shipped genome, as disclosed at the end of this subsection.}
\label{tab:periter}
\begin{tabular}{@{}llrrr@{}}
\toprule
Execution & Width & $K=16$ & $K=64$ & $K=256$ \\
\midrule
Device loop  & float64                & 22.54 & 20.39 & 19.78 \\
Device loop  & float32 iterates       & 22.99 & 20.41 & 19.78 \\
Host phases  & float64                & 27.18 & 16.19 & 13.29 \\
Host phases  & float32 throughout     & 28.00 & 12.23 &  8.43 \\
\bottomrule
\end{tabular}
\end{table}

Width buys $1.00\times$ inside the device loop at $K=256$, where the measured
ratio is $1.00006$ and the section's own stability figure is $0.66$ percent, so
the third decimal is below resolution. It buys $1.576\times$ once the loop is
chunked into host phases, with the float32 host-phase cell taken from the
patched policy registry rather than the shipped genome. Execution structure is
worth $1.488\times$ at float64 and $2.345\times$ at float32 at $K=256$, and
reverses to $0.829\times$ at $K=16$, where chunking pays more in launches than
it saves. Under PDHG the float64 execution-structure gain at $K=256$ is
$1.398\times$ rather than $1.488\times$. Reporting one of these without the
other misattributes the gain.

% source: results/artifacts_20260910/fanova.json, same block, percent_variance.
% Response is log per-iteration cost; factors are execution, precision, cadence.
% This is NOT the algorithm-by-precision-by-cadence decomposition on log
% penalised SGM-1, whose algorithm share is 0.00 percent at 1e-4 and 0.34 at
% 1e-6 and which is reported in the introduction and in Section 2
% (2-related-work.tex lines 96 to 98). Do not merge the two decompositions.
A three-way decomposition of log per-iteration cost over those twelve cells
assigns $47.43$ percent of the variance to cadence, $25.96$ percent to the
execution-by-cadence interaction, $15.98$ percent to execution structure,
$3.26$ percent to execution by precision, $2.90$ percent to precision alone,
$2.43$ percent to precision by cadence and $2.04$ percent to the three-way
residual. The design is saturated at one observation per cell, so these are
shares of an exact decomposition and carry no error bar.

% source: results/finish_20260910/loop_mechanism.json, 48 rows, recomputed.
% The 48 rows are 2 sizes x 6 trip counts x 2 traced x 2 certificate settings,
% so the certificate term is a difference and there are 24 of them, not 48.
Table~\ref{tab:loopmech} settles the cause. It holds the total iteration budget
fixed at 512 and moves the number of outer trips from 32 down to 1, with the
cadence moving inversely.

\begin{table}[t]
\centering
\caption{Traced minus static cost, nanoseconds per iteration, certificate off.
The product of outer trips and cadence is 512 in every column. One measurement
per cell, no repeats.}
\label{tab:loopmech}
\begin{tabular}{@{}lrrrrrr@{}}
\toprule
Outer trips & 32 & 16 & 8 & 4 & 2 & 1 \\
\midrule
$n=128$ penalty & 11972 & 11365 & 10893 & 11014 & 10823 & 10769 \\
$n=128$ ratio   & 2.374 & 2.348 & 2.299 & 2.361 & 2.333 & 2.317 \\
$n=256$ penalty & 12943 & 12564 & 11604 & 11929 & 11202 & 11531 \\
$n=256$ ratio   & 2.091 & 2.126 & 1.989 & 2.060 & 1.938 & 1.997 \\
\bottomrule
\end{tabular}
\end{table}

The penalty is per iteration and nearly flat in the number of outer trips. It
falls $10.0$ percent at $n=128$, from $11972$ to $10769$\,ns, and $10.9$ percent
at $n=256$, from $12943$ to $11531$\,ns, while the outer count falls 32-fold,
and the static per-iteration cost itself moves only $6.2$ percent at $n=128$,
from $8714$ to $8175$\,ns, across the same 32-fold change in unroll depth. A
lost cross-iteration fusion opportunity would not be 10 percent of a 32-fold
change. The file carries one measurement per cell with no repeats, so any
residual trend is at the resolution of a single shot. Over the 24 matched
differences the certificate term, cost with the certificate on minus cost with
it off, spans $-239.90$ to $+844.56$\,ns against a loop penalty near
$11000$\,ns, so it is not a per-certificate-check cost either, which is the
explanation we previously advanced. The largest certificate term, $+844.56$\,ns,
falls at $n=128$ and 32 outer trips, the finest cadence in the file, where a
real certificate cost would appear, and is 7 percent of the loop penalty there.

% source: results/kernel_fusion_20260909/e2e/full_4090_sz{128,256,384,496}.json
% and e2e_bw/full_bw_sz{128,256,384,496}.json
End to end the static bound alone is worth $1.673\times$, $1.347\times$,
$1.291\times$ and $1.093\times$ at $n=128$, $256$, $384$ and $496$ on the 4090
and $1.228\times$, $1.382\times$, $1.401\times$ and $1.243\times$ on the
Blackwell. The 4090 benefit decays with size and the Blackwell benefit does
not.

% source: results/tv_ffi_20260906/kernelhost_4090/factorial.json
% 19.2068 / 19.7628 / 13.2732 / 12.9007 under Condat-Vu;
% 25.0212 / 19.7619 / 15.6475 / 13.7109 under PDHG.
Under Condat-Vu the custom TV kernel is at parity with XLA on this device in
both execution structures: $19.21$ against $19.76$ microseconds per iteration in
the device loop at $K=256$, an edge of $1.029\times$, and $13.27$ against the
schedule policy's $12.90$ in host phases, a loss of $0.972\times$. Under PDHG
the same kernel loses in both, $25.02$ against $19.76$, a loss of
$1.266\times$, and $15.65$ against $13.71$, a loss of $1.141\times$. The kernel
gains $1.447\times$ from chunked execution while XLA in the same artifact gains
$1.532\times$, so both gain from chunking, the gains are not equal, and
execution structure and kernel identity remain separate decisions.

% source: results/tv_ffi_20260906/px999999_4090/factorial.json (0.00576135) and
% cadence_4090_v1/factorial.json (0.01215536 baseline, 0.00789548 catalog best).
% 0.01215536 / 0.00576135 = 2.10981; 0.01215536 / 0.00789548 = 1.53953.
% These two cells were produced by patching the policy registry, not by the
% shipped genome. Label them as diagnostics, never as search output.
The best configuration measured is not expressible in the shipped genome: XLA
float32 with host phases at $K=256$ reaches $0.00576$\,s penalised SGM-1 against
the compiled float64 and $K=64$ baseline's $0.01216$\,s, a speedup of
$2.110\times$, where the best catalog cell reaches $0.00790$\,s and
$1.540\times$. That cell was produced by patching the policy registry rather
than by the shipped genome, so it is a diagnostic that motivates a schema
change and not a search result.

\subsection{Temporal kernel fusion}
\label{sec:exp:fusion}

% source: results/artifacts_20260910/kernel_admission_4090.log, 106 lines,
% header results/artifacts_20260910/kernel_admission_4090.header.
% 17 variants x 3 sizes x 2 proximal settings = 102 comparison rows, rho = 1.9
% throughout, bitwise column "no" in every row, final line PASS.
The fused kernel is numerically admitted before it is timed. The archived
admission run compares $T$ fused iterations against $T$ sequential single-step
reference launches over 102 cases, 17 built variants at three problem shapes
and both proximal settings, all at $\rho = 1.9$. The worst absolute deviation
over all 102 cases is $4.649 \times 10^{-15}$ and the run reports PASS. No case
is bitwise identical at this $\rho$. The variant names in the log carry the
thread count, so the two builds that share the label $16\times4$, at 128 and at
256 threads, are admitted separately, and all six timed tiles appear in the
log.

% source: results/kernel_fusion_20260909/e2e/full_4090_sz{128,256,384,496}.json
% source: results/kernel_fusion_20260909/e2e_bw/full_bw_sz{128,256,384,496}.json
% Totals recomputed by summing t_s over the six instances per arm.
% NOTE FOR COAUTHOR: the best-speedup and baseline rows here repeat
% Table~\ref{tab:intro-e2e} in the introduction. Keep one of the two, or drop
% the baseline rows here and keep only the worst-tile rows.
Table~\ref{tab:fusion} gives the end-to-end fusion result against the shipped
traced baseline, with the worst tile of the same sweep beside it.

\begin{table}[t]
\centering
\caption{Time to a certified $10^{-4}$ gap over six OCT scans, $K=256$,
relative to the shipped traced-bound baseline, three timing repeats per image.
Every arm certified 6 of 6 and every arm inside a cell ran the identical
per-instance iteration counts. The worst-tile identity at $n=256$ is not
stable: a 9-repeat re-measurement one day later puts $8\times4$ worst at
$1.366$ with $32\times2$ at $1.375$, so that one cell is inside the measurement
spread and the tile name is omitted.}
\label{tab:fusion}
\begin{tabular}{@{}lrrrr@{}}
\toprule
$n$ & 128 & 256 & 384 & 496 \\
\midrule
4090 baseline (s)  & 0.06647 & 0.07733 & 0.13486 & 0.22564 \\
4090 best tile     & $8\times4$ & $24\times4$ & $24\times4$ & $32\times4$ \\
4090 best speedup  & 2.629 & 1.867 & 1.634 & 1.356 \\
4090 worst tile    & $32\times2$ & not resolved & $8\times4$ & $8\times4$ \\
4090 worst speedup & 1.320 & 1.298 & 0.926 & 0.743 \\
\midrule
Blackwell baseline (s) & 0.05732 & 0.07299 & 0.13750 & 0.19964 \\
Blackwell best tile    & $8\times4$ & $24\times4$ & $32\times4$ & $16\times2$ \\
Blackwell best speedup & 1.732 & 1.450 & 1.625 & 1.149 \\
Blackwell worst speedup & 0.973 & 1.008 & 1.100 & 0.862 \\
\bottomrule
\end{tabular}
\end{table}

% source: results/kernel_fusion_20260909/e2e/full_4090_sz256.json against
% results/finish_20260910/fusion_r9_sz256.json (9 repeats, one day later).
The optimal tile moves with size and with device: the two devices agree at
$n=128$ and $n=256$ and disagree at $n=384$ and $n=496$. Tile selection is not
optional. The wrong tile is slower than not fusing at all at the two larger
sizes on the 4090, and the best-to-worst swing is $1.99\times$ at $n=128$ and
$1.82\times$ at $n=496$. The identity of the worst tile is $8\times4$ at both
large sizes and is not resolved at $n=256$, where three repeats and nine
repeats disagree about which of $8\times4$ and $32\times2$ is worst and the two
differ by 0.7 percent under the nine-repeat run.

% source: results/kernel_fusion_20260909/e2e_bw/full_bw_sz496.json
There is a cell where the correct decision is not to fuse. On the Blackwell at
$n=496$ the best tile reaches only $0.925\times$ of the plain static-loop path
and the worst reaches $0.862\times$ of the baseline, which belongs in the same
paragraph as the $1.149\times$ headline for that cell.

% source: results/kernel_fusion_20260909/e2e/full_4090_sz*.json (128 to 496),
% large/large_4090_m{2,3,4}.json and bw_day2/large_bw_m{2,3,4}.json (992 to 1984)
% CAVEAT: the three montage sizes use four images per arm, not six, and are
% built by tiling several distinct scans, because a raw scan caps an honest
% crop near 496 px. The pixel sizes 992/1488/1984 are the m2/m3/m4 montage
% convention; the args block of those three files records size 496 and n 4.
Table~\ref{tab:ushape} sweeps size to 1984 pixels. The fusion gain over the
static loop is U-shaped on the 4090: it falls monotonically from $1.571\times$
at 128 pixels to a minimum of $1.060\times$ at 992 pixels and rises to
$1.410\times$ at 1984 pixels. The Blackwell row is not monotone on either limb,
running $1.410$, $1.049$, $1.160$, $0.925$, $0.989$, $1.033$ and $1.123$, and
its minimum sits at 496 pixels. The 4090 trough is a single point and it sits
exactly at the boundary between the six-crop protocol and the four-montage
protocol, so we read the shape and not the location. Launch-overhead removal
accounts for the small end and reduced memory traffic for the large end.

\begin{table}[t]
\centering
\caption{Fusion gain over the static loop and static-loop gain over the traced
bound, by problem size. Sizes 992 and above are montages of several distinct
scans, four images per arm rather than six, so the two limbs of the 4090 row
use different instance construction.}
\label{tab:ushape}
\begin{tabular}{@{}lrrrrrrr@{}}
\toprule
Size (px) & 128 & 256 & 384 & 496 & 992 & 1488 & 1984 \\
\midrule
4090 fusion over static      & 1.571 & 1.386 & 1.266 & 1.240 & 1.060 & 1.346 & 1.410 \\
Blackwell fusion over static & 1.410 & 1.049 & 1.160 & 0.925 & 0.989 & 1.033 & 1.123 \\
4090 static over traced      & 1.673 & 1.347 & 1.291 & 1.093 & 1.016 & 0.989 & 0.993 \\
Blackwell static over traced & 1.228 & 1.382 & 1.401 & 1.243 & 1.063 & 1.060 & 1.037 \\
\bottomrule
\end{tabular}
\end{table}

The static bound decays monotonically on the 4090 and reverses past about 1000
pixels, to $0.989\times$ at 1488 and $0.993\times$ at 1984, while on the
Blackwell it does not reverse in this range. At the two largest 4090 sizes the
best recipe is the traced bound plus fusion rather than the static bound plus
fusion, so the two execution decisions are not separable at every scale.

% source: results/kernel_fusion_20260909/cadence/cad_4090_K{64,128,256,512}.json
% source: results/kernel_fusion_20260909/bw_day2/cad_bw_K{64,128,256,512}.json
% CORRECTION relative to results/kernel_fusion_20260909/REPORT.md section 6,
% which states the iteration count is 512 at every K in this sweep. It is not.
Across cadence the argmax is stable and its value is not. At $n=256$ on the
4090 the fusion gain over the static loop is $1.258\times$, $1.265\times$,
$1.336\times$ and $1.394\times$ at $K=64$, $128$, $256$ and $512$, and the best
tile is $24\times4$ at three of four cadences, losing at $K=128$ to $16\times4$
by 0.6 percent. The Blackwell selects $24\times4$ at all four with gains
$1.006\times$, $1.024\times$, $1.053\times$ and $1.062\times$. Iteration count
is not constant across this sweep, at $3712$, $3840$, $3840$ and $4608$ total
over the six images, so only the gain-over-static-loop column is a clean
per-iteration comparison and we report only that column.

\subsection{The method axis and the construction gate}
\label{sec:exp:method}

% source: results/method_axis_20260909/factorial.json
% 2944 / 682.6667 = 4.3125 exactly, at f64_K256 and tolerance 1e-4.
% SGM-1: condat_vu 0.05592390 / 0.01349855 = 4.142957;
%        pdhg      0.05635931 / 0.01343025 = 4.196444.
Scheme identity is inert on the method axis and step sizes are decisive. At
$n=256$ under \texttt{f64\_K256} and tolerance $10^{-4}$ the banked setting
$(\tau, \sigma, \rho) = (0.05, 1.0, 1.9)$ gives mean $682.7$ iterations under
both schemes and the balanced setting $(0.25, 0.25, 1.5)$ gives $2944.0$ under
both, a ratio of $4.3125$ on iteration count under both schemes, and on
penalised SGM-1 time $4.1430$ under Condat-Vu and $4.1964$ under PDHG. At
$10^{-6}$ the difference becomes feasibility: banked certifies 5 of 6 at mean
$2966.7$ iterations and balanced certifies 0 of 6, every instance running to the
5000-iteration cap, though balanced does certify 6 of 6 at $10^{-4}$.

% source: results/netflow_20260909/grid_4090.json (both error strings read
% verbatim: the Condat-Vu rows print "t*L_f/2 + t*s*||L||^2 < 1 (lhs=...)" and
% the PDHG rows print "t*s*||L||^2 = ..."; the recorded strings are 1.02267,
% 1.0304, 1.03577, 1.03781).
% source: results/structure_20260909/structure_summary_4090.json
% (raw_t025_s025_gate_lhs = 1.0226746, 1.0304043, 1.0357716, 1.0378125).
The admissible region is family dependent, and the difference is one of kind.
The same balanced parameters that are merely slow on TV are rejected before any
arithmetic runs on every min-cost flow instance: all 24 balanced netflow rows
fail with gate left-hand side $1.02267$, $1.0304$, $1.03577$ and $1.03781$ at
grid sizes 12, 16, 24 and 32 against a threshold of 1, and all 24 banked rows
run. The Condat-Vu rows record $t L_f / 2 + t s \lVert L \rVert^2$ and the PDHG
rows record $t s \lVert L \rVert^2$; the two coincide here because the min-cost
flow objective has no smooth term. The structure probe reproduces those four
values at $t = s = 0.25$ as $1.0226746$, $1.0304043$, $1.0357716$ and
$1.0378125$, against $0.51994$ to $0.52005$ across the three degree-2 regular
cells, which are admitted.

% source: results/structure_20260909/structure_summary_4090.json
% totals: runs 2106, admitted 1755, rejected 351, capped 257,
% gpu_exec_seconds 987.81.
Admission is a function of the gate product alone: over 2106 runs the gate
admits 351 of 351 settings at each of $\theta \in \{0.1, 0.3, 0.6, 0.9,
0.99\}$ and 0 of 351 at $\theta = 1.02$, admitting 1755 runs in total,
rejecting 351, capping 257 at the iteration limit and consuming 987.8 GPU
seconds. Rejected settings were never executed, so the gate is verified as
sufficient for non-divergence and never as necessary, and 257 of the admitted
runs still failed to certify inside the cap.

% source: results/structure_20260909/structure_summary_4090.json (norm bounds)
% source: results/finish_20260910/omega_ot_4090.json
% (Lnorm = 8.15999999999979, 11.5399826689641, 14.1335345897619;
%  1.02*sqrt(64) = 8.16, relative errors 2.6e-14, 3.2e-14, 1.1e-14)
How the region scales with instance size is also a property of the family.
Min-cost flow norm bounds are bounded in size, at $4.0451$ to $4.0749$ on the
grid over sizes 12 to 32 and, on the regular topology across 512, 1024 and 2048
nodes, $2.8843$ to $2.8844$ at degree 2, $3.4781$ to $3.4803$ at degree 3 and
$3.9364$ to $3.9395$ at degree 4. Optimal transport is the opposite case, and
measurement confirms its closed form: the recorded norm bound is $8.160000$,
$11.539983$ and $14.133535$ at the $32\times32$, $64\times64$ and $96\times96$
cells, which is $1.02\sqrt{m+n}$ to a relative error of $1.1 \times 10^{-14}$
to $3.2 \times 10^{-14}$, fourteen significant figures rather than bitwise
agreement, the $1.02$ being the safety inflation and the residual being
power-iteration accuracy. The admissible product is $0.0150183$, $0.0075091$
and $0.0050061$, in ratios $2.000$ and $1.500$ that match $m+n = 64$, $128$ and
$192$ exactly, so that region shrinks as $1/(m{+}n)$.

% source: results/finish_20260910/omega_mcf.json
% Lnorm 4.323183 at 2 commodities and 4.567979 at 4, both on the 16x16 grid,
% identical across the three seeds; 4.567979^2 / 4.323183^2 = 1.1165.
The corpus now contains capacitated multi-commodity flow runs, at two commodity
counts on one topology. Doubling the commodity count from 2 to 4 on the
$16\times16$ grid moves the norm bound from $4.323183$ to $4.567979$ and
shrinks the admissible product by a factor of $1.117$, not by a factor of 2. Two
points do not identify a rate, so we make no asymptotic claim about how that
family's admissible region contracts in the number of commodities.

\subsection{Predicting the step ratio from measured structure}
\label{sec:exp:omega}

% source: results/finish_20260910/omega_4090.json, omega_bw.json (native
%   omega0 over the 24 cell-and-seed groups: 6.9691 to 9.1254);
% source: results/structure_20260909/structure_summary_4090.json
%   (omega0 over the 13 probe cells: 7.5058 to 7.6733);
% source: results/finish_20260910/omega_ot_4090.json
%   (native omega0 on optimal transport: 0.4464 to 0.7928).
% sweep source scripts/diagnostics/omega_sweep.py
The step-size recipe looked unpredictable for a reason visible only once the
data are pooled. At native scaling the scale ratio $\omega_0 = \|c\| / \|b\|$
barely varies within a family: it spans $6.9691$ to $9.1254$ over the 24
cell-and-seed groups of the netflow sweep, $7.5058$ to $7.6733$ over the 13
cells of the structure probe and $0.4464$ to $0.7928$ on optimal transport. A
variable that never varied was never implicated.

% source: results/finish_20260910/omega_4090.json and omega_bw.json.
% Construction: tau = sqrt(theta*r)/||L||, sigma = sqrt(theta/r)/||L||, so the
% gate product is theta by construction and the ratio is r by construction.
% The sweep solves with PDHG at rho = 1 throughout.
The sweep moves it deliberately. Rescaling $c \mapsto \alpha c$ is exactly
$(\tau, \sigma) \mapsto (\alpha \tau, \sigma / \alpha)$, which leaves the gate
product fixed and moves $r = \tau / \sigma$ by $\alpha^2$. We sweep $\alpha$
and $r$ each over $\{0.01, 0.03, 0.1, 0.3, 1, 3, 10, 30, 100\}$, four decades,
on 8 cells and 3 seeds at $\theta = 0.99$, $K = 64$, tolerance $10^{-4}$ and a
cap of 60000 iterations, solving with PDHG at $\rho = 1$: 1944 runs, 1510
converged, over 216 instances of which 209 have at least one converged run. The
gate product $\tau \sigma \lVert L \rVert^2$ equals $0.99$ to within
$3.4 \times 10^{-16}$ in all 1944 rows and exactly in 774 of them, so the sweep
moves the ratio and never the admissible region. It is worth being plain about
what that check is and is not: $\theta = 0.99$ is a fixed input to the sweep and
$\tau$ and $\sigma$ are constructed from it, so the check verifies the
construction arithmetic and is not evidence that the gate depends on $\theta$
alone.

% source: results/finish_20260910/omega_4090.json and omega_bw.json, recomputed.
% PROTOCOL: converged rows only; a group is a (cell, seed) pair sharing a value
% of log10(r) + 2 log10(alpha) and containing at least two distinct alphas and
% at least two converged runs; a group qualifies for a window only if EVERY
% alpha in it lies inside the window. Counting singleton groups as trivially
% agreeing inflates these counts.
The exponent 2 in that rescaling identity, which is a different object from the
exponent $p$ of the rule fitted below, is a measured invariance and not a
fitted number. Group runs by
$\log_{10} r + 2 \log_{10} \alpha$ within each cell and seed, and count a group
only if it holds at least two distinct values of $\alpha$ and at least two
converged runs. A group qualifies for a window only if every $\alpha$ in it
lies inside that window. Every group whose values of $\alpha$ are all at least
1 returns identical iteration counts, 169 of 169 at $\alpha \geq 1$, 83 of 83 at
$\alpha \geq 3$ and 25 of 25 at $\alpha \geq 10$. Restricting $\alpha$ to
$[0.3, 3]$ gives agreement in 71 of 72 groups and to $[0.1, 10]$ in 230 of 237.
Over the whole grid agreement is 437 of 601, and the disagreement sits where
the objective is far smaller than the constraint data, at 12 of 120 groups
whose smallest $\alpha$ is $0.01$ and 71 of 120 whose smallest is $0.03$, which
is where the stopping test loses scale invariance.

% source: results/finish_20260910/omega_4090.json and omega_bw.json.
% THE PROTOCOL, fixed once for the whole paper and stated in words in every
% caption that reports these numbers (0-abstract, 1-introduction Table 2,
% 5-experiments Tables 4 and 6, 6-conclusion Table 2):
%   metric: certified iteration count N, not wall clock;
%   instance: a (cell, seed, alpha) triple;
%   oracle: the minimum certified N over that instance's own ratio grid;
%   capping: a run that fails to certify inside the 60000-iteration cap is
%     charged 60000 iterations, so every policy is scored on every instance;
%   rule: C/omega0^p snapped to the nearest grid ratio in log space;
%   fixed baseline: the single grid ratio minimising geometric mean regret on
%     the evaluation set itself, chosen separately for each set, which favours
%     the baseline.
% There is no committed analysis script; this was recomputed from the rows.
% On wall clock from the identical rows the same computation gives 1.626 and
% 16.92 for r = 0.01, which is also the best fixed ratio on that metric, and
% the censoring counts move to 88 at the floor and 11 at the ceiling. That is
% why the metric is named in the caption.
Table~\ref{tab:regret} scores policies against a per-instance oracle. The best
single global ratio leaves $1.851\times$ geometric-mean regret and $26.04\times$
worst case, and a rule that reads the ratio off measured structure,
$r = C / \omega_0^2$, cuts both.

\begin{table}[t]
\centering
\caption{Regret in iteration count against a per-instance oracle over the 209
instances of the nine-point ratio grid, under the protocol used in every table
of this paper that reports regret. The metric is $N$ and not wall clock; the
oracle is the minimum certified iteration count over that instance's own ratio
grid; a run that fails to certify inside the 60000-iteration cap is charged
60000 iterations, so all three policies are scored on all 209 instances. The
certified column is the count each policy solved inside the cap. Both constants
in the rule rows are fitted on these same 209 instances, so those two rows are
in sample; the constant carried onto unseen families is the separate fit of
Table~\ref{tab:transfer}.}
\label{tab:regret}
\begin{tabular}{@{}llrrr@{}}
\toprule
Policy & Constant & Geomean regret & Worst case & Certified \\
\midrule
Global constant ratio      & $r = 0.01$   & 1.851 & 26.04 & 208 \\
$r = C / \omega_0^2$, geomean-optimal & $C = 0.263$  & 1.181 & 3.87 & 209 \\
$r = C / \omega_0^2$, worst-case-optimal & $C = 0.037$ & 1.190 & 3.18 & 209 \\
\bottomrule
\end{tabular}
\end{table}

% Recomputed under the protocol above: 1.85083 / 1.18080 = 1.5675;
% 26.04167 / 3.86667 = 6.7349; 1.85083 / 1.18998 = 1.5553;
% 26.04167 / 3.17647 = 8.1984; r = 0.3 gives geomean 2.0403 and worst 6.5833,
% certifying 180 of 209, and 6.5833 / 3.17647 = 2.0725. The worst-case-optimal
% constant is flat: every C from 0.0295 to 0.0407 gives the identical worst case
% 3.17647, so only two digits of C are printed. Scoring each policy only where
% it certifies instead returns 1.8274 and 23.03 for r = 0.01 on 208 instances.
The improvement over the best fixed choice is $1.567\times$ on geometric mean
and $6.73\times$ on worst case at the geomean-optimal constant, and
$1.555\times$ and $8.20\times$ at the worst-case-optimal constant. Both of
those worst-case figures are measured against $r = 0.01$, which is the
geomean-optimal fixed ratio and not the worst-case-optimal one. The
worst-case-optimal fixed ratio is $r = 0.3$, with geometric mean $2.040$ and
worst case $6.583$, certifying 180 of 209, so the like-for-like worst-case
improvement is $2.07\times$. Scoring each policy only where it certifies,
instead of charging capped runs at the cap, moves the constant-ratio row to
$1.827\times$ and $23.03\times$ on 208 instances.

% source: results/finish_20260910/omega_4090.json and omega_bw.json
One caveat travels with this result and Section~\ref{sec:exp:transfer} takes it
up. The ratio grid censors the optimum: 91 of 209 instances have their best
ratio at the floor of $0.01$ and 11 at the ceiling of $100$, so 102 of 209
optima sit on a grid edge, counting an instance as censored when any ratio
attaining its minimum certified $N$ sits at an endpoint. A naive log-log fit of
the optimal ratio against $\omega_0$ then returns a slope of $-0.9365$ under the
geometric tie average this paper uses throughout, against $-0.9309$ with ties
broken to the smallest ratio and $-0.9421$ to the largest, none of which is the
exponent.

\subsection{Transfer of the step-ratio rule across families}
\label{sec:exp:transfer}

% source: results/finish_20260910/omega_wide_bw.json (fit set, 1404 rows)
% The functional form is not ours: r = C / omega0^2 with omega0 = ||c||/||b||
% is the primal weight of PDLP, cited in Section 2.
The functional form of this rule is not new. A step ratio set as
$C / \omega_0^2$ with $\omega_0 = \lVert c \rVert / \lVert b \rVert$ is the
primal-weight heuristic of PDLP \citep{Applegate2021Practical}, which this paper
already credits. What is measured here is how much it is worth against the best
fixed ratio and whether one constant transfers across families.

% source: results/finish_20260910/omega_wide_bw.json.
% Fitting p and C jointly on this set alone by minimising geometric-mean regret,
% over p on a 0.05 grid and C on a 0.01 dex grid, gives p = 2.00 and C = 0.9772
% at geometric-mean regret 1.18583. The argmin in C is a single grid point at
% that resolution: the neighbouring points 0.9550 and 1.0000 give 1.18894 and
% 1.19487. The earlier claim of a plateau from 0.9772 to 0.9886 does not
% reproduce and is deleted.
% Refining to p on a 0.01 grid and C on a 0.001 dex grid moves the argmin to
% p = 1.98, C = 1.040, at regret 1.18573, which is 0.008 percent lower. The
% coarser argmin is kept because that difference does not identify a rule.
% This pair, and no other, is used in the abstract, the introduction and the
% conclusion. The earlier drafts printed C = 1.10 at p = 2 (regret 1.1881) and
% C = 1.16 at p = 1.8 (regret 1.1981); both are worse on the set they claim to
% be fitted on and both are deleted.
We widen the netflow ratio grid to 13 points spanning $10^{-4}$ to $10^{2}$ and
fit $r = C / \omega_0^p$ on that set alone, by minimising geometric-mean regret
over $p$ on a $0.05$ grid and $C$ on a $0.01$ dex grid.
The fit returns $p = 2.00$ and $C = 0.977$, at geometric-mean regret $1.186$.
That pair is then carried unchanged onto three families the fit never saw, and
it is the only rule this paper reports.

% Table~\ref{tab:transfer} uses the protocol stated above Table~\ref{tab:regret},
% unchanged: charge a non-certifying run at the 60000-iteration cap, score every
% policy on every instance, and pick the fixed ratio per evaluation set.
Table~\ref{tab:transfer} is the result, under the protocol already used for
Table~\ref{tab:regret}. Policies are scored in iteration count against a
per-instance oracle, the minimum certified $N$ over that instance's own ratio
grid. An instance is a (cell, seed, $\alpha$) triple. A run that fails to
certify inside the 60000-iteration cap is charged 60000 iterations, so both
policies are scored on every instance of every set. The fixed-ratio arm is the
single grid ratio minimising geometric-mean regret on the evaluation set
itself, chosen separately for each set, which favours the baseline; its
certified count is reported alongside the rule's because the transferred rule
certifies every instance of every set.

\begin{table}[t]
\centering
\caption{The $\omega_0$ rule carried off the family it was fitted on. The rule
is $r = 0.977 / \omega_0^2$, fitted on netflow wide alone and then evaluated
unchanged. Regret is in iteration count against a per-instance oracle over that
instance's own ratio grid, with a run that fails to certify inside the
60000-iteration cap charged at the cap, so both policies score every instance of
every row. The fixed column is the single grid ratio minimising geometric-mean
regret on that set, chosen separately per set, and the certified columns are the
instances each of the two policies solved inside the cap. The pooled row is the
geometric mean over the 162 held-out instances, each family carrying its own
best fixed ratio.}
\label{tab:transfer}
\begin{tabular}{@{}lrlrrrrr@{}}
\toprule
& & & \multicolumn{2}{c}{geomean regret} & & \multicolumn{2}{c}{certified} \\
\cmidrule(lr){4-5} \cmidrule(lr){7-8}
Evaluation set & $n$ & $\omega_0$ range & fixed & rule & rule worst & rule & fixed \\
\midrule
netflow wide (fitted on) & 108 & 0.0697 to 825 & 2.779 & 1.186 & 2.12 & 108 & 100 \\
optimal transport        &  54 & 0.0054 to 79.3 & 1.428 & 1.179 & 1.56 &  54 & 54 \\
multi-commodity flow     &  54 & 0.0227 to 312 & 2.840 & 1.032 & 1.26 &  54 & 48 \\
held-out topologies      &  54 & 0.0735 to 806 & 2.514 & 1.190 & 1.83 &  54 & 44 \\
\midrule
pooled held out          & 162 & 0.0054 to 806 & 2.168 & 1.131 & 1.83 & 162 & 146 \\
\bottomrule
\end{tabular}
% source: results/finish_20260910/omega_wide_bw.json (108 instances, 13 ratios)
% source: results/finish_20260910/omega_ot_wide.json (54 instances, 13 ratios)
% source: results/finish_20260910/omega_mcf.json (54 instances, 12 ratios)
% source: results/finish_20260910/omega_holdout.json (54 instances, 13 ratios)
% Recomputed to four decimals: rule geomean 1.1858 / 1.1786 / 1.0318 / 1.1899
%   / 1.1311; rule worst 2.1250 / 1.5625 / 1.2623 / 1.8252 / 1.8252.
% Best fixed under the protocol in the caption: 2.7790 at r = 0.03 certifying
%   100 of 108; 1.4279 at r = 100 certifying 54 of 54; 2.8395 at r = 0.1
%   certifying 48 of 54; 2.5136 at r = 0.003 certifying 44 of 54; pooled 2.1681
%   certifying 146 of 162. A single ratio shared across the three held-out
%   families, drawn from the 11 ratios their grids have in common, is worse
%   still, r = 3 at 2.8204.
% Under the alternative protocol, scored only where certified with an 80 percent
%   coverage floor: 2.1513 at r = 0.0003 on 87 of 108; 1.4279 at r = 100 on
%   54 of 54; 2.1854 at r = 0.001 on 44 of 54; 1.9475 at r = 0.0001 on 45 of 54;
%   pooled per family 1.8249 on 143 of 162. One shared ratio under the same
%   80 percent floor gives r = 0.1 at 2.5039 on 130 of 162; with the floor
%   dropped, r = 30 at 1.7516 on 108 of 162.
% omega_0 endpoints to four decimals: 0.0697 to 825.2357; 0.0054 to 79.2830;
%   0.0227 to 311.5050; 0.0735 to 806.2762; 0.0054 to 806.2762. These are over
%   all seeds; computing them on one seed gives different and wrong endpoints.
\end{table}

The gain over the best fixed ratio is $2.34\times$ on the set the rule was
fitted on, $1.21\times$ on optimal transport, $2.75\times$ on multi-commodity
flow and $2.11\times$ on the held-out topologies, and $1.92\times$ pooled over
the 162 held-out instances. The comparison is protocol dependent and we give the
other reading here rather than in a second table. Under the alternative protocol
in which each policy is scored only on the instances it certifies, with the
fixed ratio restricted to those certifying at least 80 percent of the set, the
fixed baselines improve to $2.151\times$ on 87 of 108 netflow wide instances,
$1.428\times$ on 54 of 54 optimal transport, $2.185\times$ on 44 of 54
multi-commodity flow and $1.947\times$ on 45 of 54 held-out topologies, a pooled
$1.825\times$ on 143 of the 162 held-out instances. They are then measuring a
different and smaller problem on every row, which is why the tables do not use
that protocol. Dropping the coverage floor as well lets one shared ratio certify
barely two thirds of the pooled set, $r = 30$ at $1.752\times$ on 108 of 162.
The rule is untouched by the choice at $1.186$, $1.179$, $1.032$, $1.190$ and
$1.131$, because it certifies every instance of every set.

% source: results/finish_20260910/omega_wide_bw.json, omega_ot_wide.json,
% omega_mcf.json, omega_holdout.json. Native omega0 at alpha = 1:
% netflow 6.9691 to 8.2524; optimal transport 0.5419 to 0.7928;
% multi-commodity flow 2.2666 to 3.1151; held-out topologies 7.3548 to 8.0628.
The sample structure has to be stated plainly, because it bounds what transfer
means here. Each set is (cells $\times$ seeds) distinct problems replicated at
nine synthetic rescalings of $c$: netflow wide is 12 distinct problems at 9
values of $\alpha$, and each held-out set is 6 at 9. Variation in $\omega_0$
within a set is therefore variation in $\alpha$, and $c \mapsto \alpha c$ is
exactly $(\tau, \sigma) \mapsto (\alpha \tau, \sigma / \alpha)$. The native
$\omega_0$ spread inside netflow is only $6.97$ to $8.25$, a factor of $1.18$,
so the rule is being tested on rescalings and not on naturally occurring scale
variation. The part that is not a rescaling is the comparison across families:
optimal transport sits at native $\omega_0$ of $0.54$ to $0.79$ and
multi-commodity flow at $2.27$ to $3.12$, both genuinely away from netflow's
$7.5$, and the transferred constant is not refitted for either.

% source: the same four files. Refitting C per family at p = 2 by the same
% grid search: 0.977 on netflow wide, 31.62 on optimal transport, 1.380 on
% multi-commodity flow, 0.275 on the held-out topologies. 31.62 / 0.275 = 115.
One constant does not fit every family. Refitting $C$ per family at $p = 2$
gives $0.977$ on netflow wide, $31.62$ on optimal transport, $1.380$ on
multi-commodity flow and $0.275$ on the held-out topologies, a span of a factor
of 115 between the smallest and the largest. What Table~\ref{tab:transfer}
reports is the netflow constant used unchanged, and the residual gap on
optimal transport, $1.179$ against a per-family-optimal $1.026$, is the price
of that transfer.

% source: results/finish_20260910/omega_4090.json and omega_bw.json (narrow),
% omega_wide_bw.json (wide), omega_seeds8.json (8 seeds),
% omega_ultrawide_bw.json (ten decades of r, eleven alphas, 4 seeds),
% omega_fine_4090.json (finer ratio grid).
% CONVENTIONS, the same two used in 0-abstract, 1-introduction and 6-conclusion
% and named in the caption of Table~\ref{tab:exponent}:
%   censoring: an instance counts as censored when ANY ratio attaining its
%     minimum certified N sits at an endpoint of that instance's grid;
%   slope ties: the log-log fit of the argmin ratio against omega0 runs over the
%     instances with a converged run, with tied argmin ratios averaged
%     geometrically. Under the other two tie rules, ties to the smallest and
%     ties to the largest, the five slopes move by at most 0.011.
% Regret-optimal exponent by grid search over p on a 0.05 grid at the best C on
% a 0.01 dex grid, the same resolution as the fit above.
% Slopes to four decimals: -0.9365, -1.5412, -1.5781, -1.6834, -1.7380.
% Censoring to the convention above: 102 of 209 (91 floor, 11 ceiling), 33 of
% 108, 46 of 144, 15 of 80, 14 of 88.
Two estimators of the exponent disagree, and reducing censoring moves only one
of them. Table~\ref{tab:exponent} runs both over five ratio grids of increasing
width. The log-log slope of the argmin ratio against $\omega_0$ rises in
magnitude from $-0.94$ on the original grid, where 102 of 209 optima sit on a
grid edge, to $-1.74$ on the ultrawide ten-decade grid, where 14 of 88 do. The
regret-minimising exponent does not move with it: it is $2.10$, $2.00$, $1.90$,
$2.00$ and $2.00$ on the same five grids. The honest reading is that censoring
biases the slope estimator toward zero, that removing censoring moves it from
$-0.94$ toward $-1.74$ without reaching the regret estimator, and that the two
therefore bracket the exponent between about $1.7$ and $2.1$ and have not met.
That bracket is consistent with the value 2 implied by the primal weight, and at
its lower end below it; we do not claim the exponent is exactly 2, and no
estimator on any of the five grids returns $1.8$ as its argmin. The regret
profile is flat enough to say how finely the data resolve it. On the original
grid the set of exponents within 1 percent of the optimum is not even an
interval: it runs from $1.40$ to $2.90$ with an isolated point at $1.15$. On the
ultrawide grid, which is the least censored, that set is $1.80$ to $2.25$, with
$p = 1.8$ costing $0.8$ percent and $p = 2.4$ costing $2.2$ percent. That is the
regret estimator alone; the verdict this paper carries is the two-estimator
bracket above, about $1.7$ to $2.1$. The finer 19-point grid has now run and is
the fourth row of the table; it does not narrow
the bracket, returning $-1.68$ and $2.00$, between the wide and ultrawide grids
on both estimators.

\begin{table}[t]
\centering
\caption{Two estimators of the exponent in $r = C / \omega_0^p$, over five
ratio grids, under the two conventions used wherever this paper reports them.
Censoring is the count of instances for which any ratio attaining the minimum
certified $N$ sits at an endpoint of that instance's grid. The slope is a
log-log fit of the argmin ratio against $\omega_0$ over the instances with a
converged run, with tied argmin ratios averaged geometrically; the other two tie
rules move each slope by at most $0.011$. The regret exponent minimises
geometric-mean regret in iteration count, searched over $p$ on a $0.05$ grid at
the best constant on a $0.01$ dex grid for each exponent.}
\label{tab:exponent}
\begin{tabular}{@{}lrlrrr@{}}
\toprule
Grid & $n$ & Ratio range & Censored & Slope & Regret exponent \\
\midrule
original    & 209 & $10^{-2}$ to $10^{2}$  & 102 & $-0.94$ & 2.10 \\
netflow wide & 108 & $10^{-4}$ to $10^{2}$ & 33 & $-1.54$ & 2.00 \\
8 seeds     & 144 & $10^{-4}$ to $10^{3}$  & 46 & $-1.58$ & 1.90 \\
finer grid  &  80 & $10^{-5}$ to $3 \times 10^{2}$ & 15 & $-1.68$ & 2.00 \\
ultrawide   &  88 & $10^{-6}$ to $10^{4}$  & 14 & $-1.74$ & 2.00 \\
\bottomrule
\end{tabular}
% source: results/finish_20260910/omega_4090.json + omega_bw.json; omega_wide_bw.json;
% source: results/finish_20260910/omega_seeds8.json; omega_fine_4090.json; omega_ultrawide_bw.json
% Slopes with ties averaged geometrically: -0.9365, -1.5412, -1.5781, -1.6834,
%   -1.7380. Ties to the smallest ratio: -0.9309, -1.5307, -1.5716, -1.6762,
%   -1.7324. Ties to the largest: -0.9421, -1.5516, -1.5845, -1.6905, -1.7437.
% Best regret at the argmin exponent, C on a 0.01 dex grid: 1.1805, 1.1858,
%   1.1571, 1.1685, 1.0618.
% Ultrawide profile at the same resolution: p=1.8 gives 1.0702, p=2.0 gives
%   1.0618, p=2.4 gives 1.0853. On a 0.001 dex grid in C the p=2.0 value is
%   1.0574; every figure printed in the paper uses the 0.01 dex grid.
% Within 1 percent of the optimum, by grid: {1.15} together with 1.40 to 2.90;
%   1.85 to 2.15; 1.65 to 2.45; 1.75 to 2.05; 1.80 to 2.25.
% The finer grid is omega_fine_4090.json: 19 ratios from 1e-5 to 300, spaced by
%   a factor of about 2.6. It has run; no sentence in this paper says otherwise.
\end{table}

\subsection{Negative results}
\label{sec:exp:negative}

% source: results/diagnostics_20260909/roofline_4090.json (model inputs) and
% results/diagnostics_20260909/REPORT.md lines 26-29 (prose table, all four
% rows including the 0.567 fitted-launch-cost row). The four scores were
% reproduced on 2026-09-10 by running scripts/diagnostics/score_roofline.py,
% whose stdout is archived at
% results/artifacts_20260910/roofline_scores_4090.log, 26 lines: mean Spearman
% 0.355, 0.363, 0.282 and 0.567 with pick rates 1/7, 0/7, 0/7 and 1/7, and M1
% per-size correlations -0.943 at n=128 and +0.829 at 992, 1488 and 1984.
Cheap a-priori predictors fail on the task where $\omega_0$ succeeds.
Table~\ref{tab:roofline} scores four analytic models of the fusion tile against
measurement by rank correlation over 7 sizes and 6 tiles.

\begin{table}[t]
\centering
\caption{Analytic tile selection against measurement. Occupancy raises the rank
correlation from $0.355$ to $0.363$ and lowers the best-tile hit rate from 1 of
7 to 0 of 7. All four models pick the best tile at most 1 time in 7, which is
the load-bearing figure.}
\label{tab:roofline}
\begin{tabular}{@{}lrr@{}}
\toprule
Model & Mean Spearman & Picked the best tile \\
\midrule
Roofline only                        & 0.355 & 1 of 7 \\
Plus occupancy                       & 0.363 & 0 of 7 \\
Plus occupancy and wave quantisation & 0.282 & 0 of 7 \\
Plus a fitted per-launch cost of $10^{-5}$\,s & 0.567 & 1 of 7 \\
\bottomrule
\end{tabular}
\end{table}

% source: results/diagnostics_20260909/roofline_4090.json
% All four device constants verified: 746.6579 GB/s, 1.1402 TFLOP/s, ridge
% 1.5271 flop/byte, unfused intensity 0.5536, 24x4 intensity 2.7253.
The interpretable structure is a change of regime. The roofline-only model
scores $-0.943$ at $n=128$ and $+0.829$ at $n \geq 992$, so it is right only
once the kernel is bandwidth-bound and is anti-correlated when the kernel is
overhead-bound. The device constants are 746.7\,GB/s achievable float64
bandwidth, 1.140\,TFLOP/s float64 and a ridge point of 1.527 flop per byte, on
an accounting of 31 flops and 56 bytes per pixel-iteration; the unfused TV
iteration sits at 0.554 flop per byte and the $24\times4$ tile reaches 2.725,
above the ridge. Two caveats belong in the text and not in a comment: the flop
count is hand-counted from \texttt{tv\_update.cu} and not instrumented, and
occupancy is computed from assumed Ada AD102 constants rather than from a
device query.

% source: results/diagnostics_20260909/sharpness_full_4090.json, recomputed.
% Log-log Pearson over the converged rows: pooled +0.9426 (17), netflow -0.5387
% (13), regular -0.8423 (6), grid -0.1197 (7), TV +0.9825 (4).
Empirical sharpness looks predictive and is not. Over 17 converged instances the
correlation between log sharpness and log iteration count is $+0.943$ pooled,
$-0.539$ within min-cost flow over 13 instances, $-0.842$ within the regular
topology over 6, $-0.120$ within the grid topology over 7 and $+0.983$ within
TV over 4. That last figure is a four-point correlation spanning only two image
sizes and is really tracking size, as the source report's own caveat says. The
pooled figure only separates two clusters, min-cost flow at mean sharpness 92.1
and mean 10230 iterations against TV at 43.3 and 624. Within a family the sign
inverts.

% source: results/diagnostics_20260909/sharpness_full_4090.json, recomputed.
% Quantity: log of the fitted decay rate of the gap over the stated window of
% certificate checks, against log N, over the 13 converged min-cost flow rows.
% Recomputed: +0.5387, +0.2006, +0.0177, -0.2751, -0.3686, -0.6201, and
% -0.7551 for the last window when the thirteenth instance is scored on its
% truncated window instead of dropped.
A leakage-free version of the estimator has to watch most of the solve before it
becomes even moderately predictive. Estimating the gap decay rate on a fixed
index window of certificate checks and correlating its logarithm against log
iteration count over the min-cost flow instances gives $+0.539$ over the first 3
checks, $+0.201$ over 5, $+0.018$ over 10, $-0.275$ over checks 10 to 20,
$-0.369$ over 20 to 40 and $-0.620$ over 40 to 80. The sign flips halfway
through: early in the solve a faster measured decay goes with more iterations
and not fewer, and only over checks 40 to 80 does the windowed estimate acquire
the negative sign that the full-solve sharpness estimate carries within
min-cost flow, $-0.539$. That window is scored on the 12 of 13 instances
that record at least 80 certificate checks; scoring the thirteenth on its
truncated window instead gives $-0.755$. Checks are 64 iterations apart, so the
window ends at 5120 iterations against a median min-cost flow solve of 9600, and
the estimator reaches $|r| = 0.620$ only after 53 percent of the run.

% source: results/structure_20260909/structure_summary_4090.json
The one structure-dependent selector that does work is the construction gate.
It is exact, symbolic and free, and admits 1755 of 2106 settings purely as a
function of the gate product, 351 of 351 at each of $\theta \in \{0.1, 0.3,
0.6, 0.9, 0.99\}$ and 0 of 351 at $\theta = 1.02$. Rejected settings were never
executed, so the gate is verified as sufficient for non-divergence and never as
necessary, and 257 of the admitted runs still hit the iteration cap. It is a
feasibility test and not a rate test, so it tells a search which recipes are
legal and nothing about which legal recipe is fast. That is the hinge between
these negatives and the case for search.

\subsection{Measurement conditions}
\label{sec:exp:conditions}

% source: results/tv_ffi_20260906/DECOMPOSITION_REPORT.md lines 6 to 8. The
% 12-process stability run has NO saved artifact; it is prose in that report,
% for the f64_K64 cell at n=256.
% source: results/finish_20260910/fusion_r9_sz256.json and fusion_r9_sz496.json
% against results/kernel_fusion_20260909/e2e/full_4090_sz{256,496}.json.
% Recomputed drift: baseline +1.97 and +1.07 percent; static over traced
% 1.3470 -> 1.4008 (+3.99 percent) and 1.0933 -> 1.1788 (+7.81 percent);
% best tile 1.8675 -> 1.9396 and 1.3556 -> 1.4279.
Timing runs were not on exclusive machines. One dispersion figure is available
as prose only: the reference cell \texttt{f64\_K64} at $n=256$, re-measured in
12 independent processes, spans 20.297 to 20.432 microseconds per iteration, a
spread of 0.66 percent, and that run has no saved artifact. A recomputable
dispersion is available and is larger. Re-running the $n=256$ and $n=496$
fusion sweeps at 9 repeats one day later moves the traced baseline by 2.0 and
1.1 percent, the static-over-traced gain from $1.347$ to $1.401$, a change of
4.0 percent, and from $1.093$ to $1.179$, a change of 7.8 percent, and the
best-tile speedup from $1.868$ to $1.940$ and from $1.356$ to $1.428$. No
fusion figure in this section is stable in its third significant figure. The
fusion tables print three digits so that the arithmetic can be checked against
the artifacts, and only the first two should be read as measured.

% source: results/method_axis_20260909/data.json ('final_holdout_loaded': false)
% source: results/hardware_codex/REPORT.md line 7 ("The original
% pre-registration and final holdout remain untouched.")
The frozen final holdout has never been evaluated, and none of the figures in
this section touch it. The families described as held out in
Section~\ref{sec:exp:transfer} are held out from the rule fit only; they are not
the frozen paper holdout.

A vLLM server held 15.4\,GB resident on the timed 4090 at 0 percent
utilisation throughout. The Blackwell box is shared: the fusion microbenchmarks
used GPU 2 at 0 percent utilisation but not exclusively, and all four GPUs were
idle for the cross-architecture end-to-end runs.
